\chapter{Conjectures in Elementary Number Theory}

Number theory is known for its wealth of open problems that have inspired mathematicians and enthusiasts to learn more about the subject. Some of the most famous problems, such as Goldbach's conjecture and the Twin Prime conjecture have remained elusive for centuries. We state here some open problems and conjectures in elementary number theory.

There are many more open problems, specifically concerning analytic number theory and advanced branches of algebraic number theory, which require tools beyond the scope of this text in order to be formulated. For a more comprehensive overview of the open problems and conjectures in number theory, see \cite{guy2004}.

\section{Conjectures About the Infinitude of Certain Classes of Primes}

One of the oldest and most famous open problems about the infinitude of certain classes of primes is the Twin Prime Conjecture. A pair of \emph{twin primes} is a pair of primes $p$ and $q$ such that $q-p=2$. In \zcref{exercise:modular-restrictions-twin-primes} we stated some elementary congruences satisfied by pairs of twin primes.

\begin{conjecture}[Twin Prime Conjecture]
  There are infinitely many twin primes.
\end{conjecture}

Although the study of twin primes dates back to antiquity, the first publication in which the twin prime conjecture was explicity stated was by Alphonse de Polignac \cite{polignac1846}. He also speculated about a generalization of the Twin Prime Conjecture.

\begin{conjecture}[de Polignac's Conjecture]
  For every even number $n$, there are infinitely many pairs of primes $p$ and $q$ such that $q-p=n$.
\end{conjecture}

A \emph{Sophie Germain prime} is a prime $p$ such that $2p+1$ is also prime. Sophie Germain studied these primes in the context of Fermat's Last Theorem. In \zcref{exercise:sophie-germain-primes} we established an elementary congruence regarding Sophie Germain primes. In \zcref{exercise:quadratic-residues-sophie-germain-primes} we characterized Sophie Germain primes in terms of their primitive roots.

\begin{conjecture}[Infinitude of Sophie Germain primes]
  There are infinitely many Sophie Germain primes.
\end{conjecture}

Dickson's conjecture generalizes both conjectures of de Polignac and the conjectured infinitude of the Sophie Germain primes. Recall that we introduced the concept of fixed divisors in \zcref{section:fixed-divisors}, where we proved \zcrefnamed{theorem:hensels-fixed-divisor-theorem}.

\begin{conjecture}[Dickson's Conjecture]
  Consider a finite set of linear functions $f_i(x)=a_ix+b_i$ indexed by $1\leq i\leq n$, which is \emph{admissible} in the sense that their product
  \begin{equation*}
    \prod_{i=1}^n f_i(x)
  \end{equation*}
  has no fixed divisors. Then there are infinitely many positive integers $x$ such that each $f_i(x)$ is prime for all $1\leq i\leq n$.
\end{conjecture}

One can also ask about the infinitude of primes that can be presented as the value of some polynomial. The simplest such conjecture is a conjecture of Landau.

\begin{conjecture}[Landau's Conjecture]
  There are infinitely many primes of the form $n^2+1$.
\end{conjecture}

Landau's conjecture is generalized by the following conjecture Bunyakovsky.

\begin{conjecture}[Bunyakovsky's Conjecture]
  Consider an irreducible polynomial $f(x)$ with integer coefficients, such that the leading coefficient is positive. If the fixed divisor of $f(x)$ is $1$, then $f(x)$ is prime for infinitely many input values $x$.
\end{conjecture}

Schinzel's Hypothesis H is a sweeping generalization of both Dickson's conjecture and Bunyakovsky's conjecture.

\begin{conjecture}[Schinzel's Hypothesis H]
  Consider a finite set of nonconstant irreducible polynomials $f_i(x)$ with integer coefficients, such that their product
  \begin{equation*}
    \prod_{i=1}^nf_i(x)
  \end{equation*}
  has no fixed divisors. Then there are infinitely many integers $x$ such that each $f_i(x)$ is prime.
\end{conjecture}

There are many more conjectures on the infinitude of primes of a certain form, or primes satisfying a certain condition. The most famous such conjecture asserts the infinitude of the Mersenne primes.

\begin{conjecture}[Infinitude of Mersenne primes]
  There are infinitely many primes of the form $M_p=2^p-1$.
\end{conjecture}

\begin{conjecture}[Infinitude of Wilson primes]
  There are infinitely many Wilson primes; that is, there are infinitely many primes $p$ such that
  \begin{equation*}
    (p-1)!\equiv -1\pmod{p^2}.
  \end{equation*}
\end{conjecture}

\begin{conjecture}[Infinitude of Wolstenholme primes]
  There are infinitely many Wolstenholme primes; that is, there are infinitely many primes $p$ such that
  \begin{equation*}
    \binom{2p-1}{p-1} \equiv 1 \pmod{p^4}.
  \end{equation*}
\end{conjecture}

\begin{conjecture}[Infinitude of Fibonacci primes]
  There are infinitely many Fibonacci primes.
\end{conjecture}

A prime number $p$ is said to be \emph{balanced} if the gap to its preceding prime is equal to the distance to its succeeding prime.

\begin{conjecture}[Infinitude of balanced primes]
  There are infinitely many balanced primes.
\end{conjecture}

\begin{conjecture}[Infinitude of Pierpont primes]
  There are infinitely many primes of the form $2^a3^b+1$.
\end{conjecture}

\begin{conjecture}[Artin's Conjecture]
  If $a$ is a nonsquare integer distinct from $-1$, then there are infinitely many primes $p$ for which $a$ is a primitive root modulo $p$.
\end{conjecture}

\section{Conjectures About Exceptional Primes}

\begin{conjecture}[Finiteness of Fermat primes]
  There are only finitely many primes of the form $F_n=2^{2^n}+1$.
\end{conjecture}

\begin{conjecture}[Finiteness of Wieferich primes]
  There are only finitely many Wieferich primes.
\end{conjecture}

Although it is conjectured that there are only finitely many Wieferich primes, it has not yet been proven that there are infinitely many non-Wieferich primes.

\begin{conjecture}[Infinitude of non-Wieferich primes]
  There are infinitely many non-Wieferich primes.
\end{conjecture}

\begin{conjecture}[Finiteness of Wal--Sun--Sun primes]
  There are only finitely many Wal--Sun--Sun primes; that is, there are only finitely many primes $p$ such that $p^2$ divides the $(p- \legendre{p}{5})$\supth{} Fibonacci number. Here, the number $\legendre{p}{5}$ is the Legendre symbol of $p$ modulo $5$.
\end{conjecture}

\section{Conjectures About Gaps Between the Primes}

\begin{conjecture}[Legendre's Conjecture]\label{conjecture:legendre}
  For every positive integer $n$, there is at least one prime between $n^2$ and $(n+1)^2$.
\end{conjecture}

\begin{conjecture}[Andrica's Conjecture]
  If $p$ and $q$ are two consecutive primes, then $\sqrt{q}-\sqrt{p}<1$.
\end{conjecture}

\begin{conjecture}[Cram\'er's Conjecture]
  The asymptotic growth of the gap from a prime $p$ to the next prime is $\Oh(\log(p)^2)$.
\end{conjecture}

\section{Conjectures About Conditions for Primality}

\begin{conjecture}[Lehmer's Conjecture]
  If $\phi(n)\mid n-1$, then $n$ is prime.
\end{conjecture}

\begin{conjecture}[The Agoh--Giuga Conjecture]
  An integer $n>1$ is prime if and only if the congruence
  \begin{equation*}
    \sum_{k=1}^{n-1} k^{n-1} \equiv -1\pmod{n}
  \end{equation*}
  holds.
\end{conjecture}

\begin{conjecture}[The Bateman--Selfridge--Wagstaff Conjecture (The New Mersenne Prime Conjecture)]
  Let $p$ be an odd natural number. If two of the following three conditions hold, then so does the third:
  \begin{enumerate}
  \item There is an integer $k$ such that $p=2^k\pm 1$ or $p=4^k\pm 3$.
  \item The Mersenne number $M_p=2^p-1$ is prime.
  \item The number $(2^p+1)/3$ is prime.
  \end{enumerate}
\end{conjecture}

For the following conjecture, let $p_n$ be the $n$\supth{} prime number, starting with $p_0:=2$, $p_1:=3$, and so on. Define the $n$\supth{} \emph{primorial} by
\begin{equation*}
  p_n\# := \prod_{i=0}^n p_i.
\end{equation*}

\begin{conjecture}[Fortune's Conjecture]
  For each $n$, if $m_0$ is the least integer $m$ such that $p_n\#+m$ is prime, then $m_0$ is prime.
\end{conjecture}

\section{Conjectures About Prime Decompositions}

\begin{conjecture}[Square-freeness of Fermat numbers]
  Every Fermat number $F_n=2^{2^n}+1$ is square-free \cite{braywarren1967}.
\end{conjecture}

\begin{conjecture}[Square-freeness of Mersenne numbers with prime exponents]
  Every Mersenne number $M_p$ with prime exponent $p$ is square-free \cite{rotkiewicz1965,braywarren1967}.
\end{conjecture}

In \zcref{exercise:prime-divisors-feit-thompson-conjecture} we showed that for any two distinct primes $p$ and $q$, any common prime divisor $r$ of the numbers
\begin{equation*}
  \frac{p^q-1}{p-1} \qquad\text{and}\qquad \frac{q^p-1}{q-1}
\end{equation*}
satisfies $r\equiv 1\pmod{2pq}$.
  
\begin{conjecture}[The Feit--Thompson Conjecture]
  There are no two distinct primes $p$ and $q$ such that
  \begin{equation*}
    \frac{p^q-1}{p-1} \qquad\text{divides}\qquad \frac{q^p-1}{q-1}
  \end{equation*}
\end{conjecture}

\begin{conjecture}[Grimm's Conjecture]
  If the numbers $n+1,\ldots,n+k$ are all composite, then there is a set of $k$ distinct primes $p_i$ such that $p_i\mid n+i$ for each $1\leq i\leq k$.
\end{conjecture}

A \emph{Sierpi\'{n}ski number} is an odd number $k$ such that every number of the form $k\cdot 2^n+1$ is composite.

\begin{conjecture}[Sierpi\'nski's problem]
  The smallest Sierpi\'nski number is $78,557$.
\end{conjecture}

A \emph{Riesel number} is an odd number $k$ such that every number of the form $k\cdot 2^n-1$ is composite.

\begin{conjecture}[Riesel's problem]
  The smallest Riesel number is $509,203$.
\end{conjecture}

\section{Conjectures About Arithmetic Functions}

\begin{conjecture}[The nonexistence of odd perfect numbers]
  There are no odd perfect numbers.
\end{conjecture}

\begin{conjecture}[Infinitude of Amicable Numbers]
  There are infinitely many pairs of amicable numbers; that is, there are infinitely many pairs of distinct numbers such that the sum of the proper divisors of either one of them equals the other.
\end{conjecture}

A number $n$ is called \emph{solitary} if for all $m\neq n$ we have
\begin{equation*}
  \frac{\sigma(m)}{m}\neq \frac{\sigma(n)}{n}.
\end{equation*}
In other words, $n$ is solitary if $n$ is the only number $m$ such that $\frac{\sigma(m)}{m}=\frac{\sigma(n)}{n}$. For example, we have
\begin{equation*}
  \frac{\sigma(10)}{10}=\frac{1+2+5+10}{10}=\frac{9}{5}.
\end{equation*}

\begin{conjecture}[The Solitary Number Conjecture]
  The number $10$ is solitary.
\end{conjecture}

\begin{conjecture}[Carmichael's Conjecture]
  If $k=\phi(n)$ for at least one $n$, then $k=\phi(n)$ for at least two distinct input values of $n$.
\end{conjecture}

\section{Conjectures About Diophantine Problems}

\begin{conjecture}[The Erd\H{o}s--Moser Conjecture]
  The only solution to the equation
  \begin{equation*}
    1^k+2^k+\cdots + m^k=(m+1)^k
  \end{equation*}
  is $1+2=3$.
\end{conjecture}

\begin{conjecture}[Brocard's Problem]
  The only integer solutions to the equation
  \begin{equation*}
    n!+1=m^2
  \end{equation*}
  are $4!+1=5^2$, $5!+1=11^2$, and $7!+1=71^2$.
\end{conjecture}

\begin{problem}[Existence of Perfect Euler Bricks]
  An \textbf{Euler brick} is a rectangular cuboid whose edges and face diagonals all have integer lengths. Does there exist a \textbf{perfect} Euler brick; that is, an Euler brick whose space diagonal is also an integer?
\end{problem}

\begin{conjecture}[The Fermat--Catalan Conjecture]
  There are only finitely many solutions in the positive integers to the equation
  \begin{equation*}
    x^a+y^b=z^c
  \end{equation*}
  satisfying $\frac{1}{a}+\frac{1}{b}+\frac{1}{c}<1$.
\end{conjecture}

\begin{conjecture}[Beal's Conjecture]
  If $a,b,c,x,y,z> 2$ are integers satisfying the equation
  \begin{equation*}
    x^a+y^b=z^c,
  \end{equation*}
  then $\gcd(x,y,z)>1$.
\end{conjecture}

\begin{conjecture}[Goormaghtigh Conjecture]
  The only solution in integers $x>y>1$ with $m,n>2$ satisfying
  \begin{equation*}
    \frac{x^m-1}{x-1}=\frac{y^n-1}{y-1}
  \end{equation*}
  are
  \begin{equation*}
    %3,2,5,7
    \frac{5^3-1}{5-1}=\frac{2^5-1}{2-1} \qquad\text{and}\qquad
    %2,90,13,3
    \frac{90^3-1}{90-1}=\frac{2^{13}-1}{2-1}.
  \end{equation*}
  That is, the only two integers that are repunits of length $>3$ in two different bases are
  \begin{equation*}
    31=(111)_5=(11111)_2\qquad\text{and}\qquad 8191=(111)_{90}=(1111111111111)_{2}.
  \end{equation*}
\end{conjecture}

\begin{conjecture}[The $abc$-Conjecture]
  Let $a$, $b$, and $c$ be relatively prime positive integers such that $a+b=c$.
  Let $\operatorname{rad}(n)$ denote the \emph{radical} of $n$, the product of the distinct prime factors of $n$.
  For every $\epsilon > 0$, there exist only finitely many triples $(a,b,c)$ such that
  \begin{equation*}
    c > \operatorname{rad}(abc)^{1+\epsilon}.
  \end{equation*}
\end{conjecture}

\begin{conjecture}[Hall's Conjecture]
  For any $\varepsilon>0$, there exists a constant $C_\varepsilon$ such that for any two positive integers $x$ and $y$ such that $x^3\neq y^2$, we have
  \begin{equation*}
    |y^2-x^3|>C_\varepsilon x^{\frac{1}{2}-\varepsilon}.
  \end{equation*}
\end{conjecture}

\begin{conjecture}[The Casas--Alvero Conjecture]
  Let $f(x)$ be a polynomial of degree $n$ with integer coefficients. If for each $0\leq k<n$ there is an integer $r_k$ such that both $f(r_k)=0$ and $f^{(k)}(r_k)=0$, then $f(x)$ must be of the form $c(x-r)^n$.
\end{conjecture}

\section{Conjectures About Representations of Numbers}

\begin{problem}[The congruent number problem]
  Find a complete characterization of all rational numbers that appear as the area of a right triangle with rational sides.
\end{problem}

\begin{conjecture}[Goldbach's Conjecture]
  Every even number greater than $2$ can be written as the sum of two primes.
\end{conjecture}

In light of Vieta's formulas for the roots of a quadratic polynomial, we can reformulate Goldbach's conjecture as follows: \emph{For every even number $n>2$ there exists a positive integer $c$ such that both roots of the polynomial $x^2+nx+c$ are primes.}

\begin{conjecture}[Lemoine's Conjecture]
  Every odd integer greater than $5$ can be written in the form $p+2q$, where $p$ and $q$ are primes.
\end{conjecture}

\begin{conjecture}[The Erd\H{o}s--Straus Conjecture]
  Every rational number of the form $4/n$ can be written as the sum of three unit fractions: There are integers $a$, $b$, and $c$ such that
  \begin{equation*}
    \frac{4}{n}=\frac{1}{a}+\frac{1}{b}+\frac{1}{c}.
  \end{equation*}
\end{conjecture}

\begin{conjecture}[Pollock's Conjecture]
  Every natural number is the sum of five tetrahedral numbers.
\end{conjecture}

\begin{conjecture}[The Sums of Three Cubes Conjecture]
  Every natural number $n\nequiv \pm 4\pmod{9}$ can be written as a sum of three cubes.
\end{conjecture}

For the following conjecture, observe that every integer $a>1$ occurs only a finite number of times in Pascal's triangle. That is, the integer $a$ can be represented in only finitely many ways as a binomial coefficient $a=\binom{n}{k}$. This follows immediately from the observation that any occurrences of the integer $a$ in Pascal's triangle are within the first $a+1$ rows. However, David Singmaster observed that the \emph{number of times} any integer occurs in Pascal's triangle seems to be bounded \cite{singmaster1971}. Indeed, currently there are no known numbers that occur more than eight times in Pascal's triangle, and the number $3003$ is the only known number that occurs eight times.

\begin{conjecture}[Singmaster's Conjecture]
  There exists a finite number $N$ such that every integer $a>1$ occurs at most $N$ times in Pascal's triangle.
\end{conjecture}

\begin{conjecture}[Rudin's Conjecture]
  The maximum number of squares that can occur in an arithmetic progression of length $n$ is asymptotically given by $\Oh(\sqrt{n})$. Furthermore, it is conjectured that the maximum is always attained in the arithmetic progression $24k+1$ for $0\leq k\leq n-1$ \cite{rudin1960}.
\end{conjecture}

\section{Conjectures About Iterative and Dynamical Processes}

\begin{conjecture}[Collatz's Conjecture]
  Let $f$ be the map that sends $n$ to $n/2$ if $n$ is even and sends $n$ to $3n+1$ if $n$ is odd. Starting from any number, one can reach the number $1$ by iteratively applying $f$.
\end{conjecture}

The Euclid--Mullin sequence $p_n$ is defined by setting $p_0:=2$ and letting $p_{n+1}$ be the least prime factor of
\begin{equation*}
  1+\prod_{k=0}^np_k.
\end{equation*}

\begin{conjecture}[The Mullin--Shanks Conjecture]
  Every prime occurs in the Euclid--Mullin sequence.
\end{conjecture}

\begin{conjecture}[Gilbreath's Conjecture]
  Let $(g_0(n))_{n\in\mathbb{N}}$ be the sequence of primes starting with $g_0(0):=2$, $g_0(1):=3$, and so on. Define
  \begin{equation*}
    g_{n+1}(k):=|g_n(k+1)-g_n(k)|
  \end{equation*}
  so that the sequence $g_{n+1}$ consists of the absolute differences between consecutive entries in the sequence $g_n$. Then $g_n(0)=1$ for every $n>0$.
\end{conjecture}

\begin{conjecture}[The Lonely Runner Conjecture]
  Suppose there are $n$ runners who run with distinct constant speeds on a circular track of length $1$. Each runner will, at some time, be at distance $1/n$ from all the others.
\end{conjecture}

\begin{conjecture}[The Catalan--Dickson Conjecture] Define the aliquot sequence of a number $n$ by $s_0(n):=n$ and $s_{k+1}(n)=\sigma(s_k(n))-s_k(n)$. That is, $s_{k+1}(n)$ is the sum of the proper divisors of $s_k(n)$. For every $n$, the sequence $(s_k(n))_{k\in\mathbb{N}}$ is eventually periodic.
\end{conjecture}

\begin{conjecture}[The 196 Problem (Lychrel Numbers)]
  For any integer $n$ written in base $10$, let $n'$ be the number written in base $10$ obtained by reversing the digits of $n$. Define the sequence $(a_n)$ by $a_0:=196$, and $a_{n+1}=a_n+a_n'$. Then none of the numbers in the sequence $(a_n)$ is palindromic.
\end{conjecture}

%%% Local Variables:
%%% mode: latex
%%% TeX-master: "elementary-number-theory"
%%% TeX-engine: luatex
%%% End:
